\documentclass[12pt,a4paper]{amsart}

\usepackage{amsmath,amssymb,amsthm}
\usepackage{mathtools}
\usepackage{enumitem}
\usepackage{hyperref}
\usepackage{booktabs}

%%─── Theorem environments ────────────────────────────────────────────
\newtheorem{theorem}{Theorem}[section]
\newtheorem{lemma}[theorem]{Lemma}
\newtheorem{proposition}[theorem]{Proposition}
\newtheorem{corollary}[theorem]{Corollary}
\newtheorem{conjecture}[theorem]{Conjecture}
\theoremstyle{definition}
\newtheorem{definition}[theorem]{Definition}
\newtheorem{problem}[theorem]{Open Problem}
\theoremstyle{remark}
\newtheorem{remark}[theorem]{Remark}
\newtheorem*{remark*}{Remark}

%%─── Macros ──────────────────────────────────────────────────────────
\newcommand{\ghat}{\widehat{g}}
\newcommand{\Phihat}{\widehat{\Phi}}
\newcommand{\Ai}{\mathrm{Ai}}
\newcommand{\Bi}{\mathrm{Bi}}
\newcommand{\ustar}{u_*}
\newcommand{\usaddle}{u^*}
\newcommand{\Fc}{F_c}
\newcommand{\norm}[1]{\|#1\|}
\newcommand{\abs}[1]{|#1|}
\newcommand{\ip}[2]{\langle #1,\, #2\rangle}
\DeclareMathOperator{\sgn}{sgn}
\newcommand{\OO}{\mathcal{O}}
\newcommand{\Zcross}{\mathcal{Z}_{\mathrm{cross}}}
\newcommand{\tstar}{t_*}

%%─── AMS metadata ────────────────────────────────────────────────────
\subjclass[2020]{42C05, 34E20, 33E10, 47B06, 11M06}
\keywords{Prolate spheroidal wave functions, Mellin transform,
effective bandwidth, crossover zone, composite asymptotic,
Langer--Olver, stationary phase}

\begin{document}

\title[Crossover Zone and $c^{-1/2}$ Upper Bound for PSWFs]{%
Composite Asymptotics in the Crossover Zone
and a Sharp $c^{-1/2}$ Upper Bound
for Prolate Spheroidal Wave Functions}

\author{[Author]}
\address{[Address]}
\email{[Email]}

\maketitle

\begin{abstract}
We study the Fourier--Mellin transforms
$\widehat{\Phi}_n^{(c)}$ of rescaled prolate
spheroidal wave functions (PSWFs) in the
high-bandwidth regime $n \le \kappa c$.
A central obstacle, identified in Paper~VI
\cite{PaperVI}, is the \emph{crossover zone}
\[
\Zcross := \bigl\{t : c^{-2/3} \le |t - t_*(c,n)|
\le c^{-1/3}\bigr\},
\]
in which neither the Airy-profile approximation
nor the algebraic WKB bound applies uniformly.
We resolve this obstacle by constructing a
\emph{uniform composite Airy--WKB expansion}
valid throughout $\Zcross$:
\[
\widehat{\Phi}_n^{(c)}(t)
= c^{-1/2}\,\mathcal{A}_n^{(c)}\!
  \bigl(c^{1/2}(t-\tstar)\bigr)
\,\bigl(1 + O(c^{-1/6})\bigr),
\]
with $|\mathcal{A}_n^{(c)}(x)| \le
C_\kappa(1+|x|)^{-1/4}$ uniformly in
$n \le \kappa c$ (Theorem~A).
As a consequence we prove the sharp upper bound
\[
\Delta_\varepsilon(c,n)
\;\le\; C(\kappa,\varepsilon)\,c^{-1/2},
\]
improving the $c^{-2/3}$ result of Paper~VI
and matching the scaling observed numerically
in Paper~V \cite{PaperV} (Theorem~B).
We also show that global $L^2$ mass-budget
arguments cannot yield the corresponding lower
bound $\Delta_\varepsilon \ge c_1 c^{-1/2}$,
because the dominant $L^2$-mass resides in the
algebraic zone $|t - \tstar| \ge c^{-1/3}$,
far from $t_*$.
Any proof of the lower bound requires
\emph{pointwise non-cancellation control} of
$\widehat{\Phi}_n^{(c)}$ near $t_*$
(Problem~\ref{prob:nonvanish}; Paper~VIII).
\end{abstract}

%%═══════════════════════════════════════════════════════════════════
\section*{Summary of Main Results}
%%═══════════════════════════════════════════════════════════════════

\paragraph{Mass distribution (corrected).}
A key correction to Paper~VI, Remark~5.3:
the $L^2$-mass in the crossover zone is
$\asymp c^{-17/12}$, not $c^{-5/6}$ as stated there.
The corrected zone table is:

\begin{center}
\renewcommand{\arraystretch}{1.4}
\begin{tabular}{llll}
\toprule
Zone & $t$-range relative to $t_*$
& $|\widehat{\Phi}_n^{(c)}(t)|$
& $L^2$-mass \\
\midrule
Airy & $|t-t_*| \lesssim c^{-2/3}$
& $\asymp c^{-1/2}$
& $\sim c^{-5/3}$ \\
Crossover $\Zcross$
& $c^{-2/3} \le |t-t_*| \le c^{-1/3}$
& $\asymp c^{-1/2}|t-t_*|^{-1/4}$
& $\asymp c^{-17/12}$ \\
Algebraic & $|t-t_*| \ge c^{-1/3}$
& $\le C_\kappa c^{-1/4}t^{-1/4}$
& $\sim O(1)$ \textbf{(bulk)} \\
Barrier & $t \ge 2t_*$
& $\le e^{-\gamma_\kappa c^{1/2}}$
& negligible \\
\bottomrule
\end{tabular}
\end{center}

\noindent\textbf{Consequence.}
The bulk of the $L^2$-mass (the full $\tfrac{1}{2}$)
resides in the \emph{algebraic zone}.
Any lower bound on $\Delta_\varepsilon$
that relies only on global mass estimates
will overshoot to $\Delta \gtrsim c$
rather than $\Delta \gtrsim c^{-1/2}$
(Proposition~\ref{prop:lower-barrier}).

\paragraph{Chain of rigorous results.}
\[
\underbrace{\text{Composite in }\Zcross}
_{\text{Thm.~A}}
\;\Longrightarrow\;
\underbrace{|\widehat{\Phi}_n^{(c)}(t)|
\le C_\kappa c^{-1/2}(1+c^{1/2}|t-t_*|)^{-1/4}}
_{\text{Cor.~\ref{cor:uniform-decay}}}
\;\Longrightarrow\;
\underbrace{\Delta_\varepsilon \le C(\kappa,\varepsilon)\,c^{-1/2}}
_{\text{Thm.~B}}.
\]

%%─── Guide for the Referee ───────────────────────────────────────────
\section*{Guide for the Referee}

This paper has a single main thread:
establish the $c^{-1/2}$ \emph{upper bound} on $\Delta_\varepsilon$
by controlling the crossover zone.
The lower bound is not established and the reasons
why it is genuinely harder are explained precisely
in Section~\ref{sec:barrier}.
The logical structure is:
\begin{enumerate}
\item Section~\ref{sec:setup}: notation and
Paper~VI results used.
\item Section~\ref{sec:mass}: corrected mass
upper bound in $\Zcross$
(Proposition~\ref{prop:cross-mass};
proof uses only (VI.2), independent of Theorem~A),
followed by the mass partition
(Lemma~\ref{lem:mass-partition})
and the corrected zone table
(Corollary~\ref{cor:zone-table}).
\item Section~\ref{sec:composite}: construction
and rigorous justification of the composite
asymptotic (Theorem~A);
followed by the matching lower bound on the
crossover mass (Corollary~\ref{cor:cross-mass-lower}).
\item Section~\ref{sec:upper}: proof of the
$c^{-1/2}$ upper bound (Theorem~B).
\item Section~\ref{sec:barrier}: structural
analysis showing why the lower bound is
a different problem (Proposition~\ref{prop:lower-barrier}).
\item Section~\ref{sec:problems}: open problems.
\end{enumerate}

\tableofcontents

%%═══════════════════════════════════════════════════════════════════
\section{Setup and Reminders from Paper~VI}
\label{sec:setup}
%%═══════════════════════════════════════════════════════════════════

We retain all notation from Papers~V--VI.
Recall:
\begin{itemize}
\item $V_c(u) = c^2 e^{4u}
- (\chi_n^{(c)} + \tfrac{3}{4})\,e^{2u}
+ \tfrac{1}{4}$,
the potential in Schr\"odinger form
(Paper~VI, eq.~(2.2);
note the sign convention
$-\Phi'' + V_c \Phi = 0$,
so $V_c > 0$ in the classically
forbidden region $u \ll \ustar$).

\begin{remark}[Sign convention for $V_c$]
\label{rem:sign}
Throughout this paper we use the
Schr\"odinger form
$-(\Phi_n^{(c)})'' + V_c(u)\Phi_n^{(c)} = 0$
with
\[
V_c(u) = c^2 e^{4u}
- \bigl(\chi_n^{(c)} + \tfrac{3}{4}\bigr)e^{2u}
+ \tfrac{1}{4}.
\]
The turning point $\ustar$ is the unique root
of $V_c$ in $(-\infty, -\tfrac{1}{2}\log c)$,
satisfying $V_c(u) < 0$ for $u < \ustar$
(classically allowed region)
and $V_c(u) > 0$ for $u > \ustar$
(forbidden region).
The WKB momentum is
$p(u;c) = \sqrt{-V_c(u)}$
for $u < \ustar$
and $p(u;c) = \sqrt{V_c(u)}$
for $u > \ustar$.
This is consistent with
Papers~V and VI throughout.
\end{remark}
\item $\ustar = \ustar(c,n)$, unique turning point
of $V_c$ in $(-\infty, -\tfrac12\log c)$.
\item $p(u;c) = \sqrt{V_c(u)^+}$, WKB momentum.
\item $\tstar = t_*(c,n) = p(\ustar;c) \asymp c^{1/2}$,
the saddle frequency.
\item $\usaddle = \usaddle(t;c,n)$, saddle point
defined by $p(\usaddle;c) = t$, for
$0 < t < \tstar$.
\item $\zeta(u;c,n)$, Langer variable,
$\tfrac{2}{3}|\zeta|^{3/2}
= |\int_{\ustar}^u p\,ds|$.
\item $D_n^{(c)}$, normalisation constant,
$|D_n^{(c)}| \asymp c^{1/4}$
(Paper~VI, Theorem~2.1).
\end{itemize}

The following results from Paper~VI are
used throughout without reproof:

\begin{enumerate}[(VI.1)]
\item\label{VI1} \textbf{Uniform Airy approximation}
(Thm.~2.1): For all $u \in (-\infty, 0)$,
\[
\Phi_n^{(c)}(u)
= \frac{D_n^{(c)}}{p(u;c)^{1/6}}
\Ai(\zeta(u;c,n))\,\bigl(1 + O(c^{-2/3})\bigr).
\]
\item\label{VI2} \textbf{Pointwise Mellin decay}
(Thm.~3.1): For $1 \le t \le c$,
\[
|\widehat{\Phi}_n^{(c)}(t)|
\le C_\kappa\,c^{-1/4}\,t^{-1/4}.
\]
\item\label{VI3} \textbf{Log-free $L^2$ tail}
(Thm.~3.2): For $1 \le R \le c^{1/4}$,
\[
\|\widehat{\Phi}_n^{(c)}\|_{L^2(|t|>R)}^2
\le C_\kappa^2\,c^{-1/2}\,R^{-1/2}.
\]
\item\label{VI4} \textbf{Fourier--Airy profile}
(Lem.~4.2): For $|t - \tstar| \le c^{-1/3}$,
\[
\widehat{\Phi}_n^{(c)}(t)
= c^{-1/2}\,\kappa_n^{(c)}\,
\Ai\!\bigl(c^{2/3}(t-\tstar)\bigr)
\,(1 + O(c^{-1/3})),
\quad
|\kappa_n^{(c)}| \asymp 1.
\]
\item\label{VI5} \textbf{$L^2$ normalisation}
(eq.~(2.4)):
$\|\widehat{\Phi}_n^{(c)}\|_{L^2(\mathbb{R})}^2
= \tfrac{1}{2}$.
\end{enumerate}

\begin{definition}[Crossover zone]
\label{def:Zcross}
\[
\Zcross(c,n)
:= \bigl\{t \in \mathbb{R} :
c^{-2/3} \le |t - \tstar| \le c^{-1/3}\bigr\}.
\]
In the rescaled variable $x := c^{1/2}(t - \tstar)$,
this becomes $\{x : c^{-1/6} \le |x| \le c^{1/6}\}$.
\end{definition}

\begin{remark}[Why the crossover zone is difficult]
\label{rem:difficulty}
The Airy profile \eqref{VI4} is valid for
$|c^{2/3}(t-\tstar)| = O(1)$,
i.e.\ $|t-\tstar| = O(c^{-2/3})$.
The WKB formula \eqref{VI2} requires the saddle
$\usaddle$ to lie well below $\ustar$, i.e.\
$|t - \tstar| \gg c^{-1/2}$ (so that
$\ustar - \usaddle \gg c^{-1/2}$).
In the crossover zone $\Zcross$, the gap
$\ustar - \usaddle \asymp c^{-1/6}$
is too large for the Airy approximation
and too small for the pure WKB formula.
Neither approximation applies, and a
\emph{composite} treatment is necessary.
\end{remark}

\begin{remark}[Three scales and the composite bridge]
\label{rem:scales}
Three natural rescalings appear in the analysis:
\begin{itemize}
\item \textbf{Airy scale:} $\xi = c^{2/3}(t - \tstar)$,
  natural for $|t - \tstar| \lesssim c^{-2/3}$
  (Airy profile \eqref{VI4}).
\item \textbf{WKB scale:} $y = c^{1/2}(t - \tstar)$,
  natural for $|t - \tstar| \gg c^{-1/2}$
  (algebraic decay \eqref{VI2}).
\item \textbf{Composite/Langer scale:}
  $x = c^{1/2}(t - \tstar)$ (the same scaling as the WKB variable),
  but with profile function $\mathcal{A}_n^{(c)}(x)$
  that reduces to the Airy function for
  $|x| \ll c^{1/6}$ via
  $\xi = c^{-1/6}\, x\,/\,({\mu_n^{(c)}})^{2/3}$.
\end{itemize}
The Langer--Olver map (Lemma~\ref{lem:phase})
provides the explicit bridge:
$\mathcal{A}_n^{(c)}(x)
= \kappa_n^{(c)}\,\Ai\!\bigl(x/(\mu_n^{(c)})^{2/3}\bigr)$
interpolates between the Airy regime
($|x| \ll c^{1/6}$, Airy variable $\xi = c^{2/3}(t-\tstar)$)
and the WKB regime
($|x| \gg 1$, where $|\mathcal{A}| \lesssim |x|^{-1/4}$).
The composite variable $x$ is not ad hoc:
it is precisely the output of the Langer substitution
applied to the phase in Lemma~\ref{lem:phase}.
\end{remark}

%%═══════════════════════════════════════════════════════════════════
\section{Corrected Mass Upper Bound in $\Zcross$}
\label{sec:mass}
%%═══════════════════════════════════════════════════════════════════

\begin{proposition}[Upper bound on mass in $\Zcross$]
\label{prop:cross-mass}
For all $c \ge c_0(\kappa)$ and $n \le \kappa c$:
\begin{equation}
\label{eq:cross-mass-upper}
\int_{\Zcross}
|\widehat{\Phi}_n^{(c)}(t)|^2\,dt
\;\le\; C_\kappa\, c^{-17/12}.
\end{equation}
In particular, $c^{-17/12} \ll 1$,
so the crossover zone carries negligible $L^2$-mass
compared to the algebraic zone.
\end{proposition}

\begin{proof}
This proof uses only \eqref{VI2}; no appeal
to Theorem~A is made.

From \eqref{VI2}: $|\widehat{\Phi}_n^{(c)}(t)|
\le C_\kappa c^{-1/4} t^{-1/4}$.
For $t \in \Zcross$ we have
$t \asymp \tstar \asymp c^{1/2}$,
so $|\widehat{\Phi}_n^{(c)}(t)|^2
\le C_\kappa^2\, c^{-1/2}\, t^{-1/2}$.

Write $t = \tstar + r$ with
$c^{-2/3} \le |r| \le c^{-1/3}$.
Since $t \ge |r|$ we have $t^{-1/2} \le |r|^{-1/2}$.
Using the substitution $t = \tstar + c^{-1/2}x$,
$dt = c^{-1/2}\,dx$, and the refined amplitude
$|\widehat{\Phi}_n^{(c)}|^2 \le C_\kappa^2 c^{-1} |x|^{-1/2}$
(from \eqref{VI2} with $t \asymp c^{1/2}$):
\begin{align}
\int_{\Zcross} |\widehat{\Phi}_n^{(c)}(t)|^2\,dt
&\le C_\kappa^2\, c^{-1} \cdot c^{-1/2}
\int_{c^{-1/6}}^{c^{1/6}} x^{-1/2}\,dx
\notag\\
&= C_\kappa^2\, c^{-3/2}\cdot
\bigl[2x^{1/2}\bigr]_{c^{-1/6}}^{c^{1/6}}
= C_\kappa^2\,c^{-3/2}\cdot 2(c^{1/12}-c^{-1/12})
\notag\\
&\le 2C_\kappa^2\, c^{-3/2+1/12}
= 2C_\kappa^2\, c^{-17/12}.
\label{eq:cross-sharp}
\end{align}
\end{proof}

\begin{remark}[Matching lower bound deferred]
\label{rem:lower-deferred}
The matching lower bound
$\int_{\Zcross}|\widehat{\Phi}_n^{(c)}|^2\,dt
\ge c_\kappa\, c^{-17/12}$
requires a pointwise lower bound on
$|\mathcal{A}_n^{(c)}(x)|$,
which is available only after Theorem~A;
see Corollary~\ref{cor:cross-mass-lower}.
\end{remark}

\begin{lemma}[Mass partition]
\label{lem:mass-partition}
\begin{equation}
\label{eq:mass-partition}
\|\widehat{\Phi}_n^{(c)}\|_{L^2(\mathrm{alg})}^2
= \tfrac{1}{2}
- \|\widehat{\Phi}_n^{(c)}\|_{L^2(\mathrm{Airy})}^2
- \|\widehat{\Phi}_n^{(c)}\|_{L^2(\Zcross)}^2
= \tfrac{1}{2} - O(c^{-17/12}).
\end{equation}
\end{lemma}

\begin{proof}
This follows directly from \eqref{VI5}
(total $L^2$-mass $= \tfrac{1}{2}$),
Paper~VI, Cor.~4.3 (Airy mass $\sim c^{-5/3}$),
and the upper bound \eqref{eq:cross-mass-upper}
($\le C_\kappa c^{-17/12}$).
Since $c^{-5/3} \ll c^{-17/12}$
(as $\tfrac{5}{3} = \tfrac{20}{12} > \tfrac{17}{12}$),
the dominant correction is $O(c^{-17/12})$.
\end{proof}

\begin{corollary}[Corrected zone table]
\label{cor:zone-table}
The $L^2$-mass of $\widehat{\Phi}_n^{(c)}$
in each frequency zone satisfies:
\begin{align}
\label{eq:mass-Airy}
\|\widehat{\Phi}_n^{(c)}\|_{L^2(\mathrm{Airy})}^2
&\sim c^{-5/3}, \\
\label{eq:mass-cross}
\|\widehat{\Phi}_n^{(c)}\|_{L^2(\Zcross)}^2
&\asymp c^{-17/12}, \\
\label{eq:mass-alg}
\|\widehat{\Phi}_n^{(c)}\|_{L^2(\mathrm{alg})}^2
&= \tfrac{1}{2} - O(c^{-17/12}).
\end{align}
In particular, the algebraic zone carries
the full $L^2$-bulk.
\end{corollary}

\begin{proof}
Equation~\eqref{eq:mass-Airy}: Paper~VI, Cor.~4.3.

Upper bound in \eqref{eq:mass-cross}:
Proposition~\ref{prop:cross-mass}
(proved in this section using only \eqref{VI2},
independently of Theorem~A).

Lower bound in \eqref{eq:mass-cross}:
Corollary~\ref{cor:cross-mass-lower}
(Section~\ref{sec:composite}), which is proved
after Theorem~A and uses only
Theorem~A\eqref{thm:A-lower} and $|\kappa_n^{(c)}| \asymp 1$.
There is no circularity: the upper bound in
\eqref{eq:mass-cross} is used in the proof of
Theorem~A only via Lemma~\ref{lem:remainder},
which relies on \eqref{VI2} alone.
The lower bound in \eqref{eq:mass-cross}
is a consequence of Theorem~A, not an input to it.

Equation~\eqref{eq:mass-alg}:
Lemma~\ref{lem:mass-partition}.
\end{proof}

%%═══════════════════════════════════════════════════════════════════
\section{Composite Asymptotic Expansion}
\label{sec:composite}
%%═══════════════════════════════════════════════════════════════════

\subsection{Main statement (Theorem~A)}

\begin{remark}[Variable dictionary for Section~\ref{sec:composite}]
\label{rem:variables}
For the reader's convenience we collect
all rescaled variables used in this section:
\begin{center}
\renewcommand{\arraystretch}{1.4}
\begin{tabular}{lll}
\toprule
Variable & Definition & Natural range \\
\midrule
$x$ & $c^{1/2}(t - \tstar)$
    & $[c^{-1/6}, c^{1/6}]$ in $\Zcross$ \\
$s$ & $c^{1/2}(u - \ustar)$
    & integration variable \\
$x'$ & $x + O(c^{-1/4})$
     & Olver-shifted frequency \\
$s'$ & Olver image of $s$
     & $s' = s(1 + O(c^{-1/4}))$ \\
$\xi$ & $x'/(\mu_n^{(c)})^{2/3}$
      & Airy argument \\
$y$  & $(\mu_n^{(c)})^{1/3} s'$
     & Airy integration variable \\
\bottomrule
\end{tabular}
\end{center}
The three zones in the $s$-variable are:
inner $|s| \le M$,
middle $M < |s| \le c^{1/6}$,
outer $|s| > c^{1/6}$.
\end{remark}

\begin{theorem}[Composite Airy--WKB asymptotic]
\label{thm:A}
Fix $\kappa \in (0,2/\pi)$.
There exist constants $C_\kappa, c_\kappa, c_0(\kappa) > 0$
and, for each $c \ge c_0$ and $n \le \kappa c$,
a function $\mathcal{A}_n^{(c)} : \mathbb{R} \to \mathbb{C}$
such that for all $t \in \Zcross$:
\begin{equation}
\label{eq:composite-A}
\widehat{\Phi}_n^{(c)}(t)
= c^{-1/2}\,\mathcal{A}_n^{(c)}\!\bigl(c^{1/2}(t-\tstar)\bigr)
\;\bigl(1 + O(c^{-1/6})\bigr),
\end{equation}
where the $O$-constant is uniform in $n \le \kappa c$
and $t \in \Zcross$, and $\mathcal{A}_n^{(c)}$ satisfies:
\begin{enumerate}[\rm(a)]
\item\label{thm:A-upper}
\textbf{Upper bound:}
$|\mathcal{A}_n^{(c)}(x)|
\le C_\kappa\,(1+|x|)^{-1/4}$
for all $x \in \mathbb{R}$.
\item\label{thm:A-lower}
\textbf{Non-vanishing on intervals of length $O(1)$:}
For every $x_0 \in \mathbb{R}$ there exists
$x_1 \in [x_0, x_0 + L_0]$,
where $L_0 \le 5$ is a universal constant
(determined by the largest gap between consecutive
zeros of $\Ai$, cf.~\cite{Olver1974}, Ch.~11),
such that $|\mathcal{A}_n^{(c)}(x_1)| \ge c_\kappa(1+|x_0|)^{-1/4}$.
Since $\mu_n^{(c)} \in [c_\kappa, C_\kappa]$
(Lemma~\ref{lem:uniform-n}),
the rescaled Airy function
$\Ai\!\bigl(x/(\mu_n^{(c)})^{2/3}\bigr)$
has zero spacing uniformly bounded above and below
independently of $n$ and $c$,
so the non-vanishing property follows from
the standard zero distribution of $\Ai$
\textup{(\cite{Olver1974}, Ch.~11)}.
\item\label{thm:A-match-airy}
\textbf{Matching to Airy zone:}
For $|x| \le c^{1/12}$,
$\mathcal{A}_n^{(c)}(x)
= \kappa_n^{(c)}\,\Ai\!\bigl(x/(\mu_n^{(c)})^{2/3}\bigr)
\,(1 + O(c^{-1/6}))$,
consistent with \eqref{VI4},
with $|\kappa_n^{(c)}| \asymp 1$.
\item\label{thm:A-match-WKB}
\textbf{Matching to algebraic zone:}
For $|x| \ge c^{1/12}$,
$|\mathcal{A}_n^{(c)}(x)| \le C_\kappa\,|x|^{-1/4}$,
consistent with \eqref{VI2}.
\end{enumerate}
\end{theorem}

\begin{remark}[Explicit form of $\mathcal{A}_n^{(c)}$]
\label{rem:A-explicit}
The profile function is
\begin{equation}
\label{eq:A-explicit}
\mathcal{A}_n^{(c)}(x)
:= \kappa_n^{(c)}\,
\Ai\!\left(\frac{x}{(\mu_n^{(c)})^{2/3}}\right),
\qquad |\kappa_n^{(c)}| \asymp 1,
\end{equation}
where
\begin{equation}
\label{eq:mu-def}
\mu_n^{(c)} := p''(\ustar;c)\,c^{-1/2}
\;\in\; [c_\kappa, C_\kappa]
\end{equation}
is the rescaled cubic-phase coefficient
(Lemma~\ref{lem:uniform-n}).
The coefficient $\kappa_n^{(c)}$ collects
three explicit factors:
\begin{equation}
\label{eq:kappa-explicit}
\kappa_n^{(c)}
= \underbrace{2\pi}_{\text{Airy integral}}
  \cdot
  \underbrace{(\mu_n^{(c)})^{-1/3}}_{\text{Jacobian}}
  \cdot
  \underbrace{\kappa_n^{(c,\mathrm{VI})}}_{\text{Paper~VI matching}}
  \cdot (1 + O(c^{-1/6})),
\end{equation}
where:
\begin{itemize}
\item $2\pi$ comes from the Airy integral
  representation \eqref{eq:airy-integral};
\item $(\mu_n^{(c)})^{-1/3}$ is the Jacobian
  of the substitution $s' = (\mu_n^{(c)})^{-1/3}y$;
\item $\kappa_n^{(c,\mathrm{VI})}$ is the
  normalisation constant from \eqref{VI4},
  satisfying $|\kappa_n^{(c,\mathrm{VI})}| \asymp 1$.
\end{itemize}
Since $\mu_n^{(c)} \in [c_\kappa, C_\kappa]$
(Lemma~\ref{lem:uniform-n}), we conclude
$|\kappa_n^{(c)}| \asymp 1$
uniformly for all $c \ge c_0(\kappa)$,
$n \le \kappa c$.
\end{remark}

\subsection{Auxiliary lemmas}

The proof of Theorem~A rests on five lemmas,
presented in logical order.
Lemma~\ref{lem:uniform-n} (uniformity in $n$)
and Lemma~\ref{lem:amp-near} (amplitude control)
are prerequisites for
Lemma~\ref{lem:remainder} (outer tail),
Lemma~\ref{lem:intermediate} (middle zone),
and Lemma~\ref{lem:tail-ext} (tail extension);
all five feed into the main proof.
Lemma~\ref{lem:phase} (Olver normal form) is
logically prior to Lemma~\ref{lem:intermediate},
though for expository clarity it is stated last.

\begin{lemma}[Uniformity of the cubic coefficient in $n$]
\label{lem:uniform-n}
There exist constants $c_\kappa, C_\kappa > 0$,
depending only on $\kappa$, such that
\[
  \mu_n^{(c)} := p''(\ustar;c)\,c^{-1/2}
  \;\in\; [c_\kappa,\, C_\kappa]
\]
uniformly for all $n \le \kappa c$ and $c \ge c_0(\kappa)$.
\end{lemma}

\begin{proof}
At the turning point $V_c(\ustar) = 0$, so
$p(u) = \sqrt{-V_c(u)}$ satisfies
$p'(\ustar)^2 = \tfrac{1}{2}V_c''(\ustar)$
by L'H\^opital.
A direct computation using
$V_c(u) = c^2 e^{4u} - \chi_n^{(c)} e^{2u} + \tfrac{1}{4}$
and the turning-point identity
$\chi_n^{(c)} e^{2\ustar} = c^2 e^{4\ustar} + \tfrac{1}{4}$
gives $V_c''(\ustar) \asymp c^2 e^{4\ustar} \asymp c$
(since $e^{2\ustar} \asymp c^{-1/2}$ at the turning point).
Hence $p''(\ustar) \asymp c^{1/2}$,
so $\mu_n^{(c)} \in [c_\kappa, C_\kappa]$
uniformly in $n \le \kappa c$.
\end{proof}

%%═══════════════════════════════════════════
%% FIX 5: lem:amp-near — explicit s-variable range
%%═══════════════════════════════════════════
\begin{lemma}[Amplitude and derivative bounds near $\ustar$]
\label{lem:amp-near}
In the rescaled variable $s = c^{1/2}(u - \ustar)$,
define the rescaled amplitude
$a(s) := c^{-1/2}\Phi_n^{(c)}(\ustar + c^{-1/2}s)$.
Then:
\begin{enumerate}[\rm(i)]
  \item $\|a\|_{L^\infty} \le C_\kappa\,c^{-1/3}$,
        uniformly in $n \le \kappa c$.
  \item $|a'(s)| \le C_\kappa\,c^{-1/4}$,
        uniformly in $|s| \le c^{1/6}$ and $n \le \kappa c$.
\end{enumerate}
\end{lemma}

\begin{proof}
\textit{(i)} From the uniform Airy approximation
\eqref{VI1} with $|D_n^{(c)}| \asymp c^{1/4}$
and $|p(\ustar)|^{-1/6} \asymp c^{-1/12}$,
we obtain
$|\Phi_n^{(c)}(u)| = O(c^{1/4} \cdot c^{-1/12}) = O(c^{1/6})$
uniformly near $\ustar$.
After the $c^{-1/2}$ Jacobian factor in the definition of $a$,
this gives $\|a\|_{L^\infty} = O(c^{1/6-1/2}) = O(c^{-1/3})$.

\textit{(ii)} Note that the rescaled variable $s = c^{1/2}(u-\ustar)$
satisfies $|s| \le c^{1/6}$ precisely when
$|u - \ustar| \le c^{-1/3}$,
so the two ranges are equivalent.
For $|s| \le c^{1/6}$ (equivalently $|u-\ustar| \le c^{-1/3}$),
differentiating \eqref{VI1} with respect to $u$:
\[
  (\Phi_n^{(c)})'(u)
  = D_n^{(c)}\Bigl[
    -\tfrac{1}{6}\frac{p'(u)}{p(u)^{7/6}}\Ai(\zeta)
    + \frac{\zeta'(u)}{p(u)^{1/6}}\Ai'(\zeta)
  \Bigr]\bigl(1 + O(c^{-2/3})\bigr).
\]
Using the exact identity $\zeta'(u) = p(u)$
from the Langer transformation
(see \cite{Olver1974}, Ch.~11),
and the standard Airy bounds
\[
  |\Ai(\zeta)| \le C(1+|\zeta|)^{-1/4},
  \qquad
  |\Ai'(\zeta)| \le C(1+|\zeta|)^{1/4},
\]
the Langer variable satisfies
$|\zeta| \le C_\kappa c^{1/3}$,
so $(1+|\zeta|)^{1/4} \le C_\kappa c^{1/12}$.
With $|p(u)| \asymp c^{1/2}$,
$|p'(u)| \asymp c^{1/4}$,
$|\zeta'(u)| = p(u) \asymp c^{1/2}$,
and $|D_n^{(c)}| \asymp c^{1/4}$,
both terms satisfy:
\[
\text{Term 1:}\quad
\tfrac{1}{6}\frac{|p'|}{|p|^{7/6}}|D_n^{(c)}||\Ai(\zeta)|
\asymp c^{1/4+1/4-7/12} = c^{1/6},
\]
\[
\text{Term 2:}\quad
\frac{|\zeta'|}{|p|^{1/6}}|D_n^{(c)}||\Ai'(\zeta)|
\le C_\kappa\, c^{1/2} \cdot c^{-1/12} \cdot c^{1/4} \cdot c^{1/12}
= C_\kappa\, c^{3/4}.
\]
Term~2 dominates; hence
$|(\Phi_n^{(c)})'(u)| \le C_\kappa\,c^{3/4}$.
Under the rescaling
$a'(s) = c^{-1}(\Phi_n^{(c)})'(\ustar + c^{-1/2}s)$,
this gives $|a'(s)| \le C_\kappa\,c^{-1/4}$.

\begin{remark*}
The bound $c^{-1/4}$ replaces the earlier value
$c^{-1/3}$, which used the incorrect estimate
$|\Ai'(\zeta)| = O(1)$.
That estimate fails when $|\zeta| \le c^{1/3} \to \infty$.
The corrected bound $c^{-1/4}$ is still sufficient
for all subsequent applications: Lemma~\ref{lem:remainder}
requires $|a'/\varphi'| \to 0$; since
$|a'(s)| \le C_\kappa c^{-1/4}$ and
$|\varphi'(s)| \gtrsim c^{1/3}$ for $|s| \ge c^{1/6}$,
one has $|a'/\varphi'| = O(c^{-7/12}) \to 0$. \checkmark
The relative error in Theorem~A is $O(c^{-1/6})$,
dominated by the outer zone (Step~2 of the proof),
not by this bound directly.
\end{remark*}
\end{proof}

\begin{lemma}[Outer zone: tail remainder]
\label{lem:remainder}
For $t \in \Zcross$ and all $c \ge c_0(\kappa)$,
$n \le \kappa c$:
\[
  \Bigl|\int_{|s| > c^{1/6}}
  e^{i\varphi(s)}\,a(s)\,ds\Bigr|
  = O(c^{-2/3}),
\]
where $\varphi(s) = xs + \tfrac{\mu_n^{(c)}}{6}s^3$
with $x = c^{1/2}(t - \tstar)$.
Relative to the main term $c^{-1/2}$,
this contributes $O(c^{-1/6})$.
\end{lemma}

\begin{proof}
For $|s| \ge c^{1/6}$ and $|x| \le c^{1/6}$:
\[
  |\varphi'(s)| = \Bigl|x + \tfrac{\mu_n^{(c)}}{2}s^2\Bigr|
  \ge \tfrac{c_\kappa}{2}s^2 - c^{1/6}
  \ge \tfrac{c_\kappa}{4}s^2
  \gtrsim c^{1/3}.
\]
Integration by parts yields
$|\text{boundary}| \le C_\kappa c^{-1/4}\cdot c^{-1/3}
= C_\kappa c^{-7/12}$
using $|a| = O(c^{-1/3})$ and $|a'| = O(c^{-1/4})$,
and the remainder integral is also $O(c^{-2/3})$
by the estimates of Lemma~\ref{lem:amp-near}.
The overall bound is $O(c^{-2/3})$.
\end{proof}

%%═══════════════════════════════════════════
%% FIX 1: lem:intermediate — O vs o clarification
%%═══════════════════════════════════════════
\begin{lemma}[Middle zone: non-stationary phase]
\label{lem:intermediate}
Let $M > 1$ be fixed \textup{(}independent of $c$\textup{)}.
After the Olver diffeomorphism of
Lemma~\textup{\ref{lem:phase}},
let $\tilde{a}(s')$ denote the transformed amplitude
\textup{(\eqref{eq:atilde-def})}
and $\tilde\varphi(s') = x's' + \mu_n^{(c)}s'^3/6$.
Then for $t \in \Zcross$, $n \le \kappa c$,
$c \ge c_0(\kappa, M)$:
\[
  \left|\int_{M < |s'| \le c^{1/6}}
  e^{i\tilde\varphi(s')}\,\tilde{a}(s')\,ds'\right|
  \le C_\kappa\,c^{-1/2}\,M^{-1}.
\]
\end{lemma}

\begin{proof}
The transformed amplitude is
\begin{equation}
\label{eq:atilde-def}
  \tilde{a}(s') := a\bigl(s(s',x)\bigr)
  \cdot \frac{\partial s}{\partial s'},
\end{equation}
satisfying
\begin{equation}
\label{eq:atilde-bounds}
  \|\tilde{a}\|_{L^\infty} = O(c^{-1/3}),
  \qquad
  |\tilde{a}'(s')| = O(c^{-1/4}),
\end{equation}
by Lemma~\ref{lem:phase}(ii) and Lemma~\ref{lem:amp-near}.

\textit{Phase lower bound.}
Fix $M > 1$ independent of $c$.
Since $t \in \Zcross$ implies $|x| \le c^{1/6}$
and $x' = x + O(c^{-1/4})$,
we have $|x'| \le c^{1/6}(1+o(1))$ for large $c$.
There exists $c_0 = c_0(\kappa, M) < \infty$ such that
$c^{1/6} \le \tfrac{c_\kappa}{4}M^2$
for all $c \ge c_0(\kappa, M)$.
All subsequent estimates hold for
$c \ge c_0(\kappa, M)$ with $M$ held constant;
the limit $c \to \infty$ is taken with $M$ fixed.
For $|s'| \ge M$:
\[
  |\tilde\varphi'(s')| = \Bigl|x' + \tfrac{\mu_n^{(c)}}{2}s'^2\Bigr|
  \ge \tfrac{c_\kappa}{2}M^2 - c^{1/6}
  \ge \tfrac{c_\kappa}{4}M^2
\]
uniformly on the intermediate zone.

\textit{Integration by parts.}
\begin{align*}
  \int_{M < |s'| \le c^{1/6}}
  e^{i\tilde\varphi}\tilde{a}\,ds'
  &=
  \left[\frac{e^{i\tilde\varphi}\tilde{a}}
             {i\tilde\varphi'}\right]_{M}^{c^{1/6}}
  - \int_{M}^{c^{1/6}}
    e^{i\tilde\varphi}
    \frac{d}{ds'}\!\left(\frac{\tilde{a}}{\tilde\varphi'}\right)
    ds'.
\end{align*}
Boundary at $|s'|=M$:
$|\tilde{a}/\tilde\varphi'| = O(c^{-1/3}M^{-2})$.
Boundary at $|s'|=c^{1/6}$:
$|\tilde\varphi'| \gtrsim c^{1/3}$,
so $|\tilde{a}/\tilde\varphi'| = O(c^{-2/3})$.
Remainder integral: $O(c^{-1/6}M^{-2})$.

The dominant boundary term is $O(c^{-1/3}M^{-2})$
and the remainder integral is $O(c^{-1/6}M^{-2})$.
For fixed $M$ (independent of $c$), the
remainder integral dominates asymptotically,
giving total $O(c^{-1/6}M^{-2})$.
Since $c^{-1/6}M^{-2} \cdot c^{1/2} = c^{1/3}M^{-2} \to \infty$
with $M$ fixed, this is not $o(c^{-1/2})$ in absolute terms;
rather, $O(c^{-1/6}M^{-2}) \le C_\kappa c^{-1/2} M^{-1}$
holds for all $c \ge c_0(\kappa, M)$
(choosing $c_0$ so that $c^{-1/6} \le M^{-1}c^{-1/2}/C$,
which is the statement $c_0 \ge (CM)^6$).
This gives the stated bound $C_\kappa c^{-1/2} M^{-1}$,
which is a convenient upper envelope
valid for sufficiently large $c$
(depending on $M$, which is held fixed).
\end{proof}

%%═══════════════════════════════════════════
%% FIX 2: lem:tail-ext — M = c^{1/12}, structural fix
%%═══════════════════════════════════════════
\begin{lemma}[Tail extension to $\mathbb{R}$]
\label{lem:tail-ext}
Let $M = M(c) := c^{1/12}$.
For $t \in \Zcross$, $n \le \kappa c$,
$c \ge c_0(\kappa)$:
\[
\left|
\int_{|s'| > M}
e^{i(x's' + \mu_n^{(c)}s'^3/6)}\,
\tilde{a}(s')\,ds'
\right|
= O(c^{-1/3} M^{-2})
= O(c^{-1/2}),
\]
uniformly in $n \le \kappa c$ and $t \in \Zcross$.
The error is of the same order as the main term
$c^{-1/2}$, and contributes only a
constant-factor perturbation to the
leading coefficient of
$\widehat{\Phi}_n^{(c)}(t)$,
absorbed into the implicit constant
of Theorem~A.

\begin{remark}[Why $M$ must grow with $c$]
\label{rem:M-growth}
For any \emph{fixed} $M > 1$ (independent of $c$),
one has
$c^{-1/3}M^{-2}/c^{-1/2} = c^{1/6}M^{-2} \to \infty$,
so the tail error would dominate the main term
and the extension would be unjustified.
The choice $M = c^{1/12}$ is the unique
scaling that balances the tail error
against $c^{-1/2}$:
\[
  c^{-1/3} \cdot (c^{1/12})^{-2}
  = c^{-1/3} \cdot c^{-1/6}
  = c^{-1/2}.
\]
This is consistent with the Olver normal form
(Lemma~\ref{lem:phase}), which is valid on
$|s|, |x| \le M$ for any $M$ growing
sufficiently slowly; $M = c^{1/12}$ satisfies
$c^{-1/2}M^4 = c^{-1/6} \to 0$,
so the quartic remainder \eqref{eq:R4-bound}
remains negligible on the inner zone
$|s| \le M = c^{1/12}$.
\end{remark}
\end{lemma}

\begin{proof}
For $|s'| > M$ and $|x'| \le c^{1/6}$,
the phase lower bound
$|\tilde\varphi'(s')| \ge \tfrac{c_\kappa}{4}M^2$
holds by the same argument as
Lemma~\ref{lem:intermediate}
(valid for $c \ge c_0(\kappa)$ with $M = c^{1/12}$,
since $c^{1/6} \le \tfrac{c_\kappa}{4}M^2 = \tfrac{c_\kappa}{4}c^{1/6}$
holds for $c_\kappa \ge 4$, which is ensured
by choosing $c_0(\kappa)$ large enough).
One integration by parts gives:
\[
\left|\int_{|s'|>M}
e^{i\tilde\varphi}\tilde{a}\,ds'\right|
\le \frac{2|\tilde{a}(M)|}{|\tilde\varphi'(M)|}
+ \int_{|s'|>M}
\left|\frac{d}{ds'}\!\left(
\frac{\tilde{a}}{\tilde\varphi'}\right)\right|ds'.
\]
The boundary term is
$\le C_\kappa c^{-1/3} \cdot \tfrac{4}{c_\kappa}M^{-2}
= O(c^{-1/3}M^{-2})$
using $|\tilde{a}| = O(c^{-1/3})$.
The integral term is also $O(c^{-1/3}M^{-2})$
by the same estimates.
With $M = c^{1/12}$:
$c^{-1/3}M^{-2} = c^{-1/3} \cdot c^{-1/6} = c^{-1/2}$,
confirming the stated bound.
\end{proof}

\begin{lemma}[Phase expansion and Olver normal form]
\label{lem:phase}
Write $t = \tstar + c^{-1/2}x$,
$u = \ustar + c^{-1/2}s$.
The rescaled oscillatory phase is
\begin{equation}
\label{eq:phase-expand}
  \Psi(s,x) := c^{1/2}[\phi_+(u;t,c) - \phi_0]
  = xs + \tfrac{\alpha_2}{2}s^2
    + \tfrac{\mu_n^{(c)}}{6}s^3
    + R_{\ge 4}(s,x),
\end{equation}
where $\phi_0 = S(\ustar;c) + \tstar\ustar$,
\[
  \alpha_2 = p'(\ustar)\,c^{-1/2} = O(c^{-1/4}),
\]
and for $|s|, |x| \le c^{1/6}$:
\begin{equation}
\label{eq:R4-bound}
  |R_{\ge 4}(s,x)| \le C_\kappa\,c^{-1/2}\,|s|^4.
\end{equation}
For $M = M(c) = c^{1/12}$ there exists
$c_0 = c_0(\kappa)$ such that for
$c \ge c_0$, $n \le \kappa c$,
there is a smooth diffeomorphism
$s \mapsto s'(s,x)$ defined on $|s|, |x| \le M$,
with $M = c^{1/12}$ growing sufficiently slowly
that the quartic remainder satisfies
$c^{-1/2}M^4 = c^{-1/6} \to 0$
(so the Olver reduction remains valid):
\begin{enumerate}[\rm(i)]
  \item \textbf{Normal form on the inner zone:}
        On $|s|, |x| \le M$,
        \[
          \Psi(s,x) = x'\,s'
          + \tfrac{\mu_n^{(c)}}{6}s'^3
          + O(c^{-1/2}M^4),
        \]
        uniformly,
        where $x' = x + O(\alpha_2) = x + O(c^{-1/4})$.
  \item \textbf{Jacobian:}
        $\partial s'/\partial s = 1 + O(\alpha_2)
        = 1 + O(c^{-1/4})$,
        so $ds' = (1 + O(c^{-1/4}))\,ds$.
  \item \textbf{Uniformity in $n$:}
        The diffeomorphism depends only on
        $\mu_n^{(c)} \in [c_\kappa, C_\kappa]$
        \textup{(Lemma~\ref{lem:uniform-n})};
        all constants depend only on $\kappa$.
\end{enumerate}

%%═══════════════════════════════════════════
%% FIX 2 (cont.): rem:M-fixed → admissible growth rate
%%═══════════════════════════════════════════
\begin{remark}[The admissible growth rate of $M$]
\label{rem:M-fixed}
The parameter $M$ must grow slowly enough
that the Olver normal form remains valid
on $|s| \le M$, i.e.\ the quartic remainder
$c^{-1/2}M^4 \to 0$ is required.
This gives the constraint $M \ll c^{1/8}$.
The choice $M = c^{1/12}$ satisfies this:
$c^{-1/2}M^4 = c^{-1/6} \to 0$. \checkmark

At the other extreme, $M = c^{1/3}$ would give
$c^{-1/2}M^4 = c^{5/6} \to \infty$,
invalidating the Olver reduction entirely.
The window of valid choices is
$1 \ll M \ll c^{1/8}$;
the choice $M = c^{1/12}$ sits comfortably
in the interior of this window and
simultaneously balances the tail-extension
error at $O(c^{-1/2})$
(Lemma~\ref{lem:tail-ext}, Remark~\ref{rem:M-growth}).
\end{remark}

\begin{remark}[Uniformity of the Olver shift in $x$]
\label{rem:olver-uniform}
The quadratic term $\tfrac{\alpha_2}{2}s^2$ is absorbed
into the Olver diffeomorphism. On $|s| \le M$,
its contribution to the phase is
$\alpha_2 s^2 = O(c^{-1/4})$ uniformly in
$x \in [-c^{1/6}, c^{1/6}]$,
since $\alpha_2 = O(c^{-1/4})$ and $M = c^{1/12}$.
The resulting frequency shift satisfies
$x' = x + O(c^{-1/4})$, with relative size
$(x'-x)/x = O(c^{-5/12}) \to 0$,
confirming that the Olver reduction is uniform in $x$.
\end{remark}
\end{lemma}

\begin{proof}
\textit{Taylor expansion.}
Expanding $S(u;c) = \int_{\ustar}^u p(v;c)\,dv$
around $\ustar$ and substituting
$u = \ustar + c^{-1/2}s$,
the linear term cancels by
$p(\ustar;c) = \tstar$ (saddle condition),
yielding \eqref{eq:phase-expand}--\eqref{eq:R4-bound}.

\textit{Olver reduction.}
On $|s|, |x| \le M = c^{1/12}$,
the quartic remainder satisfies
$c^{-1/2}M^4 = c^{-1/6} \to 0$,
and the coefficient $\alpha_2/2 = O(c^{-1/4}) \to 0$,
so the phase is a smooth small perturbation
of the pure cubic $xs + \mu_n^{(c)}s^3/6$.
Olver's normal-form theorem
(\cite{Olver1974}, Ch.~11, Theorem~11.1)
applies on this \emph{slowly growing} domain
(since the remainder $\to 0$ uniformly)
and provides the diffeomorphism satisfying~(i)--(iii).
\end{proof}

%%═══════════════════════════════════════════════════════════════════
\begin{proof}[Proof of Theorem~A]
We decompose the Fourier--Mellin transform.
Write $x = c^{1/2}(t - \tstar)$,
$a(s) = c^{-1/2}\Phi_n^{(c)}(\ustar + c^{-1/2}s)$,
so that
\[
  \widehat{\Phi}_n^{(c)}(t)
  = c^{-1/2}\int_{\mathbb{R}}
    e^{i\Psi(s,x)}\,a(s)\,ds.
\]

\medskip
\noindent\textbf{Step~1: Three-zone decomposition.}
Set $M = c^{1/12}$ (as in Lemmas~\ref{lem:phase}
and \ref{lem:tail-ext}).
Split:
\[
  \widehat{\Phi}_n^{(c)}(t)
  = c^{-1/2}\bigl[
    \mathcal{I}_{\mathrm{inn}}
    + \mathcal{I}_{\mathrm{mid}}
    + \mathcal{I}_{\mathrm{out}}
  \bigr],
\]
where the three integrals are over
$|s| \le M$, $M < |s| \le c^{1/6}$,
and $|s| > c^{1/6}$, respectively.

\medskip
\noindent\textbf{Step~2: Outer zone.}
By Lemma~\ref{lem:remainder}:
\[
  c^{-1/2}|\mathcal{I}_{\mathrm{out}}|
  = O(c^{-2/3}).
\]
This is $O(c^{-1/6})$ relative to the main term $c^{-1/2}$.

\medskip
\noindent\textbf{Step~3: Middle zone.}
By Lemma~\ref{lem:intermediate}
(applied with the current $M = c^{1/12}$,
noting $c_0(\kappa, M) = c_0(\kappa)$ for this $M$):
\[
  c^{-1/2}|\mathcal{I}_{\mathrm{mid}}|
  \le C_\kappa c^{-1/2} \cdot c^{-1/2} M^{-1}
  = C_\kappa c^{-1} \cdot c^{-1/12}
  = C_\kappa c^{-13/12}.
\]
This is $O(c^{-7/12})$ relative to $c^{-1/2}$,
hence negligible.

\medskip
\noindent\textbf{Step~4: Inner zone and Airy evaluation.}
On $|s| \le M = c^{1/12}$,
apply the Olver normal form of Lemma~\ref{lem:phase}.
Writing $s'$ for the transformed variable:
\[
  \mathcal{I}_{\mathrm{inn}}
  = \kappa_n^{(c)}
    \int_{|s'| \le M}
    e^{i(x's' + \mu_n^{(c)}s'^3/6)}\,ds'
    \cdot(1 + O(c^{-1/4})).
\]
By Lemma~\ref{lem:tail-ext} with $M = c^{1/12}$,
extending to $\mathbb{R}$ introduces error
$O(c^{-1/3}M^{-2}) = O(c^{-1/2})$,
which is a constant-factor perturbation
of the leading term and is absorbed into
the implicit constant of \eqref{eq:composite-A}.

%%═══════════════════════════════════════════
%% FIX 3: explicit Lipschitz / Ai' bound
%%═══════════════════════════════════════════
Applying the Airy integral representation
\textup{(\cite{Olver1974}, Ch.~11, eq.~(11.10.6))}:
\begin{equation}
\label{eq:airy-integral}
  \int_{\mathbb{R}}
  e^{i(\xi y + y^3/3)}\,dy
  = 2\pi\,\Ai(\xi).
\end{equation}
Setting $s' = (\mu_n^{(c)})^{-1/3}y$ gives
$\xi = x'/(\mu_n^{(c)})^{2/3}$, hence:
\[
  \int_{\mathbb{R}}
  e^{i(x's' + \mu_n^{(c)}s'^3/6)}\,ds'
  = 2\pi\,(\mu_n^{(c)})^{-1/3}\,
    \Ai\!\Bigl(\frac{x'}{(\mu_n^{(c)})^{2/3}}\Bigr).
\]

Combining with $x' = x + O(c^{-1/4})$
and the standard derivative bound
$|\Ai'(\xi)| \le C(1+|\xi|)^{1/4}$
(\cite{Olver1974}, Ch.~11),
which gives
\[
  \Bigl|\Ai\!\Bigl(\frac{x'}{(\mu_n^{(c)})^{2/3}}\Bigr)
  - \Ai\!\Bigl(\frac{x}{(\mu_n^{(c)})^{2/3}}\Bigr)\Bigr|
  \le C(1+|x|)^{1/4} \cdot
  \frac{|x'-x|}{(\mu_n^{(c)})^{2/3}}
  = O(c^{-1/4})(1+|x|)^{1/4},
\]
the replacement $x' \to x$
introduces relative error $O(c^{-1/4})$:
\[
  \mathcal{I}_{\mathrm{inn}}
  = 2\pi\,(\mu_n^{(c)})^{-1/3}\,\kappa_n^{(c)}\,
    \Ai\!\Bigl(\frac{x}{(\mu_n^{(c)})^{2/3}}\Bigr)
  + O(c^{-1/2}M^{-1}).
\]

\medskip
\noindent\textbf{Step~5: Combining all zones.}
\[
  \widehat{\Phi}_n^{(c)}(t)
  = c^{-1/2}\,\mathcal{A}_n^{(c)}(x)
  \cdot\bigl(1 + O(c^{-1/6})\bigr),
\]
where $\mathcal{A}_n^{(c)}(x)
= \kappa_n^{(c)}\,\Ai(x/(\mu_n^{(c)})^{2/3})$
and the relative error $O(c^{-1/6})$ is uniform
in $n \le \kappa c$ and $t \in \Zcross$.
This gives \eqref{eq:composite-A}.

Properties~(a)--(d) follow from:
\begin{itemize}
  \item[(a)] Standard Airy bound
        $|\Ai(\xi)| \le C(1+|\xi|)^{-1/4}$
        and $\mu_n^{(c)} \in [c_\kappa, C_\kappa]$.
  \item[(b)] Non-vanishing of $\Ai$:
        on every interval of length $L_0 \le 5$,
        $\Ai$ does not vanish identically,
        giving the lower bound
        $c_\kappa(1+|x_0|)^{-1/4}$
        (\cite{Olver1974}, Ch.~11);
        the uniform zero spacing follows from
        $\mu_n^{(c)} \in [c_\kappa, C_\kappa]$
        as stated in Theorem~A\eqref{thm:A-lower}.
  \item[(c),(d)] Matching with \eqref{VI4}
        at $x = 0$ (Remark~\ref{rem:A-explicit})
        and asymptotic WKB for $|x| \ge c^{1/12}$.
\end{itemize}
\end{proof}

\begin{remark}[Relative error convention]
\label{rem:relative-error}
All error bounds in Theorem~A and its proof are
stated in two ways:
\begin{itemize}
\item \emph{Absolute:} the bound on the contribution
  to $\widehat{\Phi}_n^{(c)}(t)$ itself
  (e.g.\ $O(c^{-2/3})$).
\item \emph{Relative:} the bound divided by the
  main-term scale $c^{-1/2}$
  (e.g.\ $O(c^{-2/3})/c^{-1/2} = O(c^{-1/6})$).
\end{itemize}
The error contributions, relative to the main term
$c^{-1/2} \cdot \mathcal{I}_{\mathrm{inn}} \asymp c^{-1/2}$, are:
\[
\frac{|\mathcal{I}_{\mathrm{out}}|}{|\mathcal{I}_{\mathrm{inn}}|}
+ \frac{|\mathcal{I}_{\mathrm{mid}}|}{|\mathcal{I}_{\mathrm{inn}}|}
+ \text{(tail ext.)}
+ \text{(Jacobian)}
= O(c^{-1/6}) + O(c^{-7/12}) + O(1) + O(c^{-1/4}).
\]
The tail-extension error $O(1)$ relative to $c^{-1/2}$
contributes only a constant-factor perturbation
to the leading coefficient $\kappa_n^{(c)}$
and is absorbed into the implicit constant of Theorem~A.
The dominant \emph{multiplicative} relative error
is $O(c^{-1/6})$ (outer zone), confirming
the $1 + O(c^{-1/6})$ in \eqref{eq:composite-A}.
\end{remark}

\begin{corollary}[Lower bound on mass in $\Zcross$]
\label{cor:cross-mass-lower}
For all $c \ge c_0(\kappa)$, $n \le \kappa c$:
\[
\int_{\Zcross}|\widehat{\Phi}_n^{(c)}(t)|^2\,dt
\;\ge\; c_\kappa\,c^{-17/12}.
\]
Together with Proposition~\ref{prop:cross-mass}
this gives $\asymp c^{-17/12}$.
\end{corollary}

\begin{proof}
By Theorem~A\eqref{thm:A-lower},
for every $x_0 \in [c^{-1/6},c^{1/6}]$
there is $x_1 \in [x_0,x_0+L_0]$ with
$|\mathcal{A}_n^{(c)}(x_1)| \ge c_\kappa|x_0|^{-1/4}$.
Hence via \eqref{eq:composite-A}:
\[
\int_{\Zcross}|\widehat{\Phi}|^2\,dt
\ge c_\kappa^2 c^{-3/2}
\int_{c^{-1/6}}^{c^{1/6}} x^{-1/2}\,dx
= c_\kappa^2 c^{-3/2}\cdot 2(c^{1/12}-c^{-1/12})
\ge c_\kappa^2 c^{-17/12}.
\]
\end{proof}

%%═══════════════════════════════════════════════════════════════════
\section{The $c^{-1/2}$ Upper Bound}
\label{sec:upper}
%%═══════════════════════════════════════════════════════════════════

\begin{corollary}[Uniform decay in $\Zcross$]
\label{cor:uniform-decay}
For $t \in \Zcross$, $c \ge c_0$, $n \le \kappa c$:
\begin{equation}
\label{eq:uniform-decay}
|\widehat{\Phi}_n^{(c)}(t)|
\le C_\kappa c^{-1/2}(1+c^{1/2}|t-\tstar|)^{-1/4}.
\end{equation}
\end{corollary}

\begin{proof}
Immediate from Theorem~A\eqref{thm:A-upper}
and \eqref{eq:composite-A},
since $x = c^{1/2}(t-\tstar)$ gives
$(1+|x|)^{-1/4} = (1+c^{1/2}|t-\tstar|)^{-1/4}$.
\end{proof}

\begin{theorem}[Upper bound on peak width, Theorem~B]
\label{thm:B}
For every $\kappa\in(0,2/\pi)$ and $\varepsilon\in(0,1)$
there exists $C(\kappa,\varepsilon)>0$ such that
for all $c \ge c_0(\kappa,\varepsilon)$, $n \le \kappa c$:
\begin{equation}
\label{eq:thm-B}
\Delta_\varepsilon(c,n) \le C(\kappa,\varepsilon)\,c^{-1/2}.
\end{equation}
\end{theorem}

\begin{proof}
We estimate the $L^2$-mass of $\widehat{\Phi}_n^{(c)}$
outside $[t_* - \Delta, t_* + \Delta]$
for $\Delta = K c^{-1/2}$, $K = K(\kappa,\varepsilon)$
chosen below.

\textbf{Step~1: Decomposition.}
Write
$\|\widehat{\Phi}_n^{(c)}\|_{L^2([t_*-\Delta, t_*+\Delta]^c)}^2
= J_1 + J_2 + J_3$,
where $J_1$ integrates over $\Zcross \cap \{|t-t_*|>\Delta\}$,
$J_2$ over the algebraic zone
$c^{-1/3} < |t-t_*| \le t_*$,
and $J_3$ over $|t-t_*| > t_*$.
Since $\Delta = Kc^{-1/2} \gg c^{-2/3}$ for large $c$,
the Airy zone lies inside $[t_*-\Delta, t_*+\Delta]$
and contributes nothing to $J_1$.

\textbf{Step~2: Bound on $J_1$.}
$J_1 \le \|\widehat{\Phi}_n^{(c)}\|^2_{L^2(\Zcross)}
\le C_\kappa c^{-17/12}$
by Proposition~\ref{prop:cross-mass}.
For $c \ge c_1(\kappa,\varepsilon)$ we have
$C_\kappa c^{-17/12} \le \varepsilon^2/6$.

\textbf{Step~3: Bound on $J_3$.}
For $t > 2t_*$, the barrier estimate
(Paper~V, Thm.~3.1) gives exponential decay;
for $t < 0$, $\widehat{\Phi}_n^{(c)} = 0$.
Hence $J_3 \le e^{-\gamma_\kappa c^{1/2}} \le \varepsilon^2/6$
for $c \ge c_2(\kappa,\varepsilon)$.

%%─── FIX 4: Step 4 — korrigierter J_2-Beweis ─────────────────────
\textbf{Step~4: Bound on $J_2$ (corrected).}
We split $J_2 = J_2^- + J_2^+$, where
$J_2^-$ integrates over $\{0 < t < t_* - \Delta\}$
and $J_2^+$ integrates over $\{t_* + \Delta < t \le 2t_*\}$.
(The region $t > 2t_*$ is absorbed into $J_3$.)

\textit{Bound on $J_2^-$.}
Since $0 < t \le t_*$ throughout this region,
\eqref{VI2} gives
$|\widehat{\Phi}_n^{(c)}(t)|^2 \le C_\kappa^2 c^{-1/2} t^{-1/2}$.
Integrating directly:
\[
J_2^-
\le C_\kappa^2 c^{-1/2}
\int_0^{t_*} t^{-1/2}\,dt
= 2C_\kappa^2 c^{-1/2} \cdot t_*^{1/2}
\asymp 2C_\kappa^2 c^{-1/4}.
\]
In particular, $J_2^- = O(c^{-1/4}) \to 0$
as $c \to \infty$, independently of $K$.

\textit{Bound on $J_2^+$.}
We apply \eqref{VI2} directly (not \eqref{VI3},
whose hypothesis $R \le c^{1/4}$ fails here
since $R = t_* + \Delta \asymp c^{1/2} \gg c^{1/4}$).
For $t \in [t_* + \Delta,\, 2t_*]$ we have
$t \ge t_* \asymp c^{1/2}$,
so $t^{-1/2} \le C c^{-1/4}$.
Hence:
\begin{align*}
J_2^+
&\le C_\kappa^2\, c^{-1/2}
\int_{t_*+\Delta}^{2t_*} t^{-1/2}\,dt \\
&\le C_\kappa^2\, c^{-1/2} \cdot C\, c^{-1/4}
\cdot (2t_* - t_* - \Delta) \\
&\le C_\kappa^2\, c^{-3/4} \cdot t_*
\;\asymp\; C_\kappa^2\, c^{-1/4}.
\end{align*}
In particular, $J_2^+ = O(c^{-1/4}) \to 0$,
independently of $K$.

\begin{remark*}
The use of \eqref{VI3} in the previous version
was erroneous: \eqref{VI3} requires $R \le c^{1/4}$,
but here $R = t_* + \Delta \asymp c^{1/2} \gg c^{1/4}$.
The corrected argument uses only \eqref{VI2},
which is valid for all $1 \le t \le c$.
The conclusion $J_2^+ = O(c^{-1/4})$
is the same as before, so Theorem~B is unaffected.
\end{remark*}

\textit{Combined.}
$J_2 = J_2^- + J_2^+ = O(c^{-1/4}) \to 0$.

\textbf{Conclusion.}
Choose $K$ so that $J_1 \le \varepsilon^2/6$,
which requires $c^{-17/12} \le \varepsilon^2/(6C_\kappa)$,
hence $K$ is determined by $c_1(\kappa,\varepsilon)$ alone.
Then for $c \ge \max(c_1,c_2,c_3)(\kappa,\varepsilon)$:
\[
J_1 + J_2 + J_3
\le \frac{\varepsilon^2}{6} + \frac{\varepsilon^2}{3}
+ \frac{\varepsilon^2}{6}
= \frac{2\varepsilon^2}{3} < \varepsilon^2.
\]
This shows $\Delta_\varepsilon(c,n) \le Kc^{-1/2}
= C(\kappa,\varepsilon) c^{-1/2}$.
\end{proof}

%%─── FIX 4 (cont.): rem:explicit-const ──────────────────────────
\begin{remark}[Explicit constant]
\label{rem:explicit-const}
In the corrected proof, $J_2$ tends to zero
independently of $K$, so $K$ is determined
only by the condition $J_1 \le \varepsilon^2/6$,
i.e.\ $C_\kappa c^{-17/12} \le \varepsilon^2/6$.
This gives
\[
C(\kappa,\varepsilon)
\;\sim\; \frac{6C_\kappa}{\varepsilon^2},
\]
a significant improvement over the original
$O(\varepsilon^{-4})$ bound arising from the
erroneous $K$-dependence in the previous Step~4.
\end{remark}

%%═══════════════════════════════════════════════════════════════════
\section{Why the Lower Bound is a Different Problem}
\label{sec:barrier}
%%═══════════════════════════════════════════════════════════════════

\begin{proposition}[Global mass arguments
cannot give the $c^{-1/2}$ lower bound]
\label{prop:lower-barrier}
Let $\mathcal{M}$ be any argument that uses
only the following two inputs:
\begin{enumerate}[\rm(a)]
\item the total $L^2$-normalisation
$\|\widehat{\Phi}_n^{(c)}\|_{L^2}^2
= \tfrac{1}{2}$,
\item pointwise upper bounds of the form
$|\widehat{\Phi}_n^{(c)}(t)|
\le f(t,c)$ for some explicit $f$.
\end{enumerate}
Then $\mathcal{M}$ cannot establish
$\Delta_\varepsilon(c,n) \ge c_1\,c^{-1/2}$
for any $c_1 > 0$.
The best lower bound that $\mathcal{M}$
can produce is
$\Delta_\varepsilon(c,n) \ge c_1({\kappa,\varepsilon})\,c^{-1}$.
\end{proposition}

\begin{proof}
A lower bound on $\Delta_\varepsilon$
via a mass argument takes the form:
\[
2\Delta \cdot f(\Delta, c)^2
< (1-\varepsilon^2) \cdot \tfrac{1}{2},
\]
giving $\Delta \ge c_1/f(\Delta,c)^2$.

\textit{Case 1: Near-peak bound.}
Using the sharpest available bound near $t_*$,
$|f| \le C_\kappa c^{-1/2}$:
\[
\int_{|t-t_*|\le\Delta}|\widehat{\Phi}|^2\,dt
\le 2\Delta \cdot C_\kappa^2 c^{-1}
= 2KC_\kappa^2 c^{-3/2} \to 0
\]
as $c \to \infty$ for any fixed $K$.
For this mass to account for the full $\tfrac12$,
one would need $\Delta \asymp c$ --- trivially useless.

\textit{Case 2: Crossover bound.}
Using $|f| \le C_\kappa c^{-1/2}|t-t_*|^{-1/4}$
at $|t-t_*|=\Delta$ in the mass inequality:
\[
2\Delta \cdot C_\kappa^2 c^{-1}\Delta^{-1/2} < \tfrac12
\;\Longrightarrow\;
\Delta^{1/2} < \frac{1}{4C_\kappa^2 c^{-1}}
\;\Longrightarrow\;
\Delta < C c^{-2/3}.
\]
This argument yields only an upper constraint
on admissible $\Delta$, not a lower bound,
hence cannot prove $\Delta_\varepsilon \ge c^{-1/2}$.

\textit{Structural obstruction.}
The dominant $L^2$-mass ($\tfrac{1}{2}$, bulk)
resides in the algebraic zone $|t - t_*| \ge c^{-1/3}$,
far from $t_*$.
Any method that does not distinguish
pointwise non-cancellation from cancellation
in this zone cannot see the $c^{-1/2}$ scale.
The best achievable bound via $\mathcal{M}$ is
$\Delta \ge c_1 c^{-1}$,
weaker than $c^{-1/2}$ by a factor of $c^{1/2}$.
\end{proof}

\begin{remark}[Structural obstruction]
\label{rem:structural}
Proposition~\ref{prop:lower-barrier}
identifies the precise obstruction:
the lower bound requires knowing that
$|\widehat{\Phi}_n^{(c)}(t)|$ does
\emph{not} cancel on a set of measure
$\gtrsim c^{-1/2}$ near $t_*$.
Global mass arguments cannot detect
local non-cancellation.
The required non-cancellation statement
is Problem~\ref{prob:nonvanish}.
\end{remark}

\begin{remark}[Conceptual significance of the barrier]
\label{rem:conceptual}
Proposition~\ref{prop:lower-barrier} shows that
the $c^{-1/2}$ scaling is inherently a \emph{local}
phenomenon near $t_*$.
It cannot be detected by global energy estimates,
but requires pointwise control of oscillatory
non-cancellation.
This distinction mirrors similar phenomena in
harmonic analysis, where the concentration of a
function on a small set is invisible to
$L^2$-norm information alone.
The required non-cancellation statement is
Problem~\ref{prob:nonvanish}.
\end{remark}

\begin{remark}[What would be needed]
\label{rem:lower-needed}
The correct lower bound requires
$|\widehat{\Phi}_n^{(c)}(t)| \ge c_\kappa c^{-1/4}t^{-1/4}$
on a set of measure $\gtrsim c^{-1/2}$
in the algebraic zone near $\tstar$;
see Problem~\ref{prob:nonvanish}.
\end{remark}

%%═══════════════════════════════════════════════════════════════════
\section{Open Problems}
\label{sec:problems}
%%═══════════════════════════════════════════════════════════════════

\begin{problem}[Non-cancellation lower bound]
\label{prob:nonvanish}
Prove that there exist $c_1(\kappa)>0$
and $E_c \subset \{|t-\tstar|\le c^{-1/2}\}$
with $|E_c| \ge c_1 c^{-1/2}$ such that
$|\widehat{\Phi}_n^{(c)}(t)| \ge c_1 c^{-1/2}$
for all $t\in E_c$, $n\le\kappa c$, $c\ge c_0$.
This is the central missing ingredient for
$\Delta_\varepsilon \ge c_1 c^{-1/2}$
(Paper~VIII).
\end{problem}

\begin{problem}[Sharp constant in Theorem~B]
\label{prob:sharp-const}
Determine the sharp constant
\[
C^*(\kappa,\varepsilon)
:= \limsup_{c\to\infty}
c^{1/2}\Delta_\varepsilon(c,\lfloor\kappa c\rfloor).
\]
The explicit bound $C(\kappa,\varepsilon)
\sim \frac{6C_\kappa}{\varepsilon^2}$
from Remark~\ref{rem:explicit-const}
is almost certainly not sharp.
\end{problem}

\begin{problem}[Uniformity in $n$]
\label{prob:uniformity}
Sharpen the $n$-dependence of
$\mu_n^{(c)}$ and $\kappa_n^{(c)}$
to obtain optimal constants in
Theorem~A and Theorem~B.
In particular, determine whether
$\mu_n^{(c)}$ is monotone in $n$
for fixed $c$.
\end{problem}

\begin{problem}[Extension beyond $n \le \kappa c$]
\label{prob:extension}
Extend Theorem~A and Theorem~B to the
superextremal regime $n > \kappa c$,
where the saddle-point structure changes
qualitatively.
\end{problem}

%%═══════════════════════════════════════════════════════════════════
\section*{Acknowledgements}
%%═══════════════════════════════════════════════════════════════════

[Acknowledgements to be added.]

\begin{thebibliography}{99}

\bibitem{PaperV}
[Author],
\textit{Effective Bandwidth and Peak-Width
Scaling for Prolate Spheroidal Wave Functions:
Numerical Evidence and Conjectures},
Paper~V of this series, preprint, 2025.

\bibitem{PaperVI}
[Author],
\textit{Pointwise Mellin Decay and
$L^2$ Tail Estimates for Prolate Spheroidal
Wave Functions},
Paper~VI of this series, preprint, 2025.

\bibitem{Olver1974}
F.~W.~J. Olver,
\textit{Asymptotics and Special Functions},
Academic Press, New York, 1974.

\bibitem{Slepian1978}
D.~Slepian,
Prolate spheroidal wave functions, Fourier
analysis, and uncertainty~V,
\textit{Bell System Tech.~J.}
\textbf{57} (1978), 1371--1430.

\bibitem{LandauPollak1961}
H.~J. Landau and H.~O. Pollak,
Prolate spheroidal wave functions, Fourier
analysis and uncertainty~II,
\textit{Bell System Tech.~J.}
\textbf{40} (1961), 65--84.

\end{thebibliography}

\end{document}
